\chapter{Ideals and Quotients}
\lecture{8}{Wed 12 Feb 2025 13:27}{Ideals and Quotient Rings}
\begin{definition}[Ideal]\label{dfn:2.1}
	A subring $I$ of a ring $R$ is a (two-sided) ideal of $R$ if for all $r\in R$ and $a\in I$, $ra,ar\in I$. Equivalently, for all $r\in R$, it follows that $rI,Ir\subseteq I$.
\end{definition}

\begin{proposition}\label{prp:2.1}
	A nonempty subset $I\subseteq R$ is an ideal of $R$ if
	\begin{itemize}
		\item $a-b\in I$ for all $a,b\in I$;
		\item $ra,ar\in I$ for all $a\in I$ and $r\in R$.
	\end{itemize}
\end{proposition}
\begin{proof}
	trivial.
\end{proof}
\begin{definition}[Principal Ideal]\label{dfn:2.2}
	Let $R$ be a commutative ring, and let $a\in R$. Define
	\[
		\left\langle a \right\rangle \coloneqq \left\{ ra\mid r\in R \right\} 
	.\]
	We call $\left\langle a \right\rangle $ the \emph{principal ideal} generated by $a$.
\end{definition}
\begin{proposition}\label{prp:2.2}
	Principal ideals are ideals.
\end{proposition}
\begin{proof}
	Let $ra,sa\in \left\langle a \right\rangle $. Then $ra-sa=(r-s)a\in \left\langle a \right\rangle $ since $r-s\in R$. By definition, $ra\in \left\langle a \right\rangle $ for all $r\in R$, so by proposition \ref{prp:2.1}, $\left\langle a \right\rangle $ is an ideal.
\end{proof}

Consider the very important example of $\mathbb{R} [x]$. Consider $I$ as the set of polynomials with a constant term equal to 0. That is, $I=\left\langle x \right\rangle $. We can thus find that $I=x\mathbb{R} [x]$ (we can factor out $x$ from every element of $\left\langle x \right\rangle $).

We can also consider ideals with multiple generators.

\begin{definition}[Generated Ideal]\label{dfn:2.3}
	Let $X\subseteq R$ be a set. Define
	\[
		\left\langle X \right\rangle \coloneqq \left\{ \sum_{i=0}^{k} r_ix_i\mid \left\{ x_i \right\}_{i=0}^k \subseteq X,\left\{ r_i \right\}_{i=0}^k \subseteq R\right\} 
	.\]
	That is, $\left\langle X \right\rangle $ is the collection of all $R$-linear combinations.
\end{definition}

The set $\left\langle X \right\rangle $ is trivially shown to be an ideal. We can define $R[x,y]\coloneqq R[x][y]$. I don't know why we're talking about this. Now consider cosets. We already know a lot about cosets from group theory; for example, $a-b\in I$ if and only if $a+I=b+I$. Let's apply this to the example $\mathbb{R} [x]$ together with the ideal $\left\langle x \right\rangle $.

Suppose that $f(x),g(x)\in\mathbb{R} [x]$ such that $f(x)+\left\langle x \right\rangle =g(x)+\left\langle x \right\rangle $. Note that we must have $f(x)-g(x)\in I$, and thus $f(0)-g(0)=0$. Equivalently, we have that $x$ divides $f(x)-g(x)$, so if we define $h(x)=f(x)-g(x)$, then $h(x)=xq(x)$ for some $q\in\mathbb{R} [x]$, and thus $h(0)=0$. Two ways of arriving at this fact. This implies that $f(0)=g(0)$, and thus $f,g$ have the same constant terms, as $f(0)=a_0$ clearly. We thus find that $f(x)+\left\langle x \right\rangle =g(x)+\left\langle x \right\rangle $ if and only if $f,g$ have the same constant term.
